\documentclass[10pt]{article}
\usepackage[margin=2.2cm]{geometry}
\usepackage{graphicx}
\usepackage{amsmath}
\usepackage{booktabs}
\usepackage[hidelinks]{hyperref}
\usepackage{xcolor}
\setlength{\parskip}{4pt}\setlength{\parindent}{0pt}

\title{\vspace{-1.2cm}Neuron mechanisms: how $d4$-dominant neurons compute degree-4 parity\\
\large \texttt{wd3\_k0\_N2048}, step 500k --- post-nonlinearity units, uniform evaluation data}
\author{}\date{September 2026}

\begin{document}\maketitle\vspace{-1.4cm}

\section*{Setup}
For each neuron we analyze the post-nonlinearity activation $a=\phi(z)$, with
$\phi=\mathrm{relu}(\tanh)$ and $z=(Vh)_n$, measured immediately after the
nonlinearity and before the residual addition. We call a neuron \emph{$d4$-dominant}
when, among the 15 binary-tree target modes, its largest Walsh coefficient is a
four-bit parity $\chi_{abcd}$ (with $q=(x_a,x_b,x_c,x_d)$ denoting the corresponding
tree quad). There are 1022 active $d4$-dominant neurons in total: 207/229/285/307 in
blocks 0--3, excluding dead units. All coefficients are estimated on $2^{16}$ uniform
samples. For each neuron, both $z$ and $a$ are decomposed over the complete 16-mode
Walsh basis of the corresponding 4-bit cube. In both cases, this 4-bit cube explains
$>99\%$ of the variance, meaning that these neurons are, to very good approximation,
functions only of those four input bits.

\section*{Mechanism: majority gates on raw bits, not composition of pairs}
For the exemplar neuron (block~2, neuron~1543, quad $d4\_0$), the pre-activation is
approximately a \emph{balanced sum of four singleton modes},
\[
z \;\approx\; \gamma\,( x_a - x_b - x_c - x_d),\qquad \gamma=0.305,
\]
with the four singleton terms accounting for $99.7\%$ of $\mathrm{Var}(z)$. In other
words, before the nonlinearity there is essentially no degree-2 or degree-4 content.
Since $z/\gamma\in\{-4,-2,0,2,4\}$, applying $a=\phi(z)$ turns this signed linear sum
into a \textbf{soft ``$\geq3$-of-4'' majority gate}: it fires most strongly,
$tanh(4\gamma)=0.84$, on the single input corner that matches the learned sign pattern
in all four coordinates, and more weakly, $\tanh(2\gamma)\approx0.56$, on the four
corners that match it in exactly three coordinates (Fig.~\ref{fig:ex}, left).

The resulting Walsh spectrum is almost completely reproduced by the
\emph{parameter-free} idealized gate $\phi(\gamma\sum_k s_k x_k)$
(Fig.~\ref{fig:ex}, middle): roughly $70\%$ of the variance lies in the four
singleton modes, roughly $20\%$ is distributed \emph{equally across all six pair
parities}, $8.8\%$ lies in the four-bit parity itself, and about $1\%$ lies in
triple modes. Crucially, the neuron is not receiving already-formed upstream $d2$
features as its main inputs. The network is therefore not implementing the
four-parity as an explicit $4=2{\cdot}2$ composition at this point; instead, this
neuron acts directly on the four raw degree-1 features.

\begin{figure}[h]\centering
\includegraphics[width=\textwidth]{fig1_exemplar_1543.pdf}
\caption{Exemplar $d4$-dominant neuron (block 2, \#1543). Left: conditional means
over the 16 cells of its quad --- a $\geq3$-of-4 majority gate. Middle: measured
Walsh spectrum (bars) vs.\ the ideal gate prediction (black ticks); colors mark the
tree-tracked modes. Right: the pre-activation is four balanced singletons.}
\label{fig:ex}
\end{figure}

\section*{The observed coverage $\approx 1/6$, explained}
Our tree-based analysis tracks only three modes on this four-bit cube:
$\{\chi_{ab},\chi_{cd},\chi_{abcd}\}$. These are the two degree-2 parities that are
actual children of the quad in the binary target tree, together with the degree-4
quad itself. The other four pair parities are real Walsh modes of the neuron, but
they are \emph{off-tree} and therefore do not count toward tree coverage.
For the ideal gate, the tracked fraction is
\[
\mathrm{cov} \;=\; \frac{a_{ab}^2+a_{cd}^2+a_{abcd}^2}{\mathrm{Var}(a)}
\;=\;0.151\quad(\text{exemplar }0.156,\ \text{population median }0.142),
\]
which matches the observed value of about $0.16$ (Fig.~\ref{fig:pop}, middle).
Thus the apparent ``$1/6$'' coverage is not the result of the neuron's variance being
symmetrically divided into six equal pieces. It is simply the fraction of the
majority-gate spectrum that happens to lie on the three modes we track in the target
tree. The remaining $\sim\!84\%$ is mostly singleton power ($\sim70\%$) plus the
four off-tree cross-pair modes
$\chi_{ac},\chi_{ad},\chi_{bc},\chi_{bd}$ ($\sim14\%$ in total). Each of these cross
pairs is about as strong as either of the two tree pairs; they are excluded only
because they are not nodes of the chosen binary tree.

\section*{Population: universal gates, plus quad-amplifiers in the last block}
This mechanism is not specific to the exemplar. Among the highest-ranked $d4$ neurons
in blocks 0--2, \textbf{100\%} satisfy the balanced-singleton input criterion
($\geq90\%$ of $\mathrm{Var}(z)$ in the four singleton modes, with balance
$\geq0.7$); across all active $d4$-dominant neurons, the fraction is 72.6\%
(Fig.~\ref{fig:pop}, left). The population contains all three possible sign-pattern
classes, 4:0/3:1/2:2 $=145/646/376$. These correspond to majority/corner detectors
centered on different regions of the four-bit cube rather than to one unique gate.

The main exception appears in the final residual block. There are 160 neurons---155
of them in block~3---with tree coverage close to 1. Their pre-activations are no longer
dominated by raw singleton features. Instead, the already-formed degree-4 parity is
itself the dominant input component: in $z$, the median variance allocation is
$84\%$ to $\chi_{abcd}$, $14\%$ to the two tree-pair modes, and only $0.5\%$ to the
singletons. These neurons are therefore better described as \textbf{pass-through
amplifiers}: blocks 0--2 create degree-4 parity through majority-gate neurons, while
their large singleton and off-tree pair components can cancel when many neurons are
combined; block~3 then receives the cleaner population-level parity feature and
re-amplifies it.

\begin{figure}[h]\centering
\includegraphics[width=\textwidth]{fig2_population.pdf}
\caption{All active $d4$-dominant neurons. Left: fraction with majority-gate input
vs.\ rank by quad power (blocks 0--2 at 1.0; block 3 lower --- see amplifiers).
Middle: tree-tracked coverage concentrates at the ideal-gate value $0.15$
($\approx$ the observed $1/6$), with the amplifier subpopulation at $\approx1$.
Right: median variance allocation --- most of a $d4$ neuron is off-tree.}
\label{fig:pop}
\end{figure}

\textbf{Takeaway.} In \texttt{wd3\_k0\_N2048}, degree-4 parity is not primarily
computed by composing the network's existing $d2$ features. Instead, most
$d4$-dominant neurons in blocks 0--2 implement soft majority/corner detectors directly
on the four underlying raw bits. Such a gate naturally contains the desired degree-4
parity, but also much larger singleton components and equal-strength degree-2
components on all six pairs, including four pairs that are not part of the target
binary tree. This is why the tree modes explain only about $1/6$ of the variance of a
typical $d4$-dominant neuron. The clean degree-4 parity is therefore best viewed as a
\emph{population-level} feature: unwanted singleton and off-tree components cancel
across neurons, while the shared parity component survives. In the final block,
dedicated amplifier neurons then receive this consolidated parity and carry it much
more directly.

\vfill\footnotesize
Artifacts: \texttt{gdrive:parity/analysis\_output/neuron\_mechanisms/wd3\_k0\_N2048/step\_00500000/}
(figures, \texttt{d4\_per\_neuron.csv}, exemplar tensors, this report; d2 corner-unit
probe in the same folder).
\end{document}
